\chapter{Jacobi Fields}
\label{chap:dc-jacobi-fields}

Faithful transcription of do Carmo, \emph{Riemannian Geometry}. Each numbered
statement keeps do Carmo's number, kind and (name); proofs keep his explicit
cross-references ("by Lemma 4.1 of Chap. 4", "from \cref{cor:dc-ch5-2-5}") so the
dependency graph is recoverable from the text.

\section{Introduction}

In this chapter we derive a first relation between the two basic concepts
introduced previously, namely geodesics and curvature. As we shall see (cf.
\cref{rem:dc-ch5-2-11}), the curvature $K(p,\sigma)$, $\sigma\subset T_p M$, determines how
fast the geodesics that start from $p$ and are tangent to $\sigma$ spread apart.
To formalize precisely this velocity of deviation of the geodesics, it is
necessary to introduce the so-called \emph{Jacobi fields}. Jacobi fields are vector
fields along geodesics, defined by means of a differential equation that arises
naturally in the study of the exponential mapping (cf. Sec. 2). Besides
furnishing the relation mentioned above, the Jacobi fields allow us to obtain a
simple characterization of the singularities of the exponential map (cf.
\cref{prop:dc-ch5-3-5}).

\section{The Jacobi equation}

Let $M$ be a Riemannian manifold and let $p\in M$. In the proof of the Gauss
Lemma we saw that if $\exp_p$ is defined at $v\in T_p M$, and if
$w\in T_v(T_p M)$, then

$$
(d\exp_p)_v w=\frac{\partial f}{\partial s}(1,0),
$$

where $f$ is a parametrized surface given by

$$
f(t,s)=\exp_p tv(s),\quad 0\le t\le 1,\quad -\varepsilon\le s\le\varepsilon,
$$

and $v(s)$ is a curve in $T_p M$ with $v(0)=v$, $v'(0)=w$.

We would like to obtain information on $|(d\exp_p)_v(w)|$, which denotes,
intuitively, the rate of spreading of the geodesics $t\mapsto\exp_p tv(s)$ that
start from $p$. This spreading is associated with the value of the sectional
curvature at $p$ with respect to the plane generated by $v$ and $w$. Also, if
$|(d\exp_p)_v(w)|=0$ with $w\neq 0$, then $v$ is a critical point of $\exp_p$. It
is convenient to study the field

$$
(d\exp_p)_{tv}(tw)=\frac{\partial f}{\partial s}(t,0)
$$

along the geodesic $\gamma(t)=\exp_p(tv)$, $0\le t\le 1$.

The basic remark is that $\frac{\partial f}{\partial s}$ satisfies a differential
equation. Since $\gamma$ is a geodesic, we have for all $(t,s)$,
$\frac{D}{\partial t}\frac{\partial f}{\partial t}=0$. Thus, from Lemma 4.1 of
Chap. 4,

$$
0=\frac{D}{\partial s}\Big(\frac{D}{\partial t}\frac{\partial f}{\partial t}\Big)
=\frac{D}{\partial t}\frac{D}{\partial s}\frac{\partial f}{\partial t}
-R\Big(\frac{\partial f}{\partial s},\frac{\partial f}{\partial t}\Big)\frac{\partial f}{\partial t}
=\frac{D}{\partial t}\frac{D}{\partial t}\frac{\partial f}{\partial s}
+R\Big(\frac{\partial f}{\partial t},\frac{\partial f}{\partial s}\Big)\frac{\partial f}{\partial t}.
$$

Putting $\frac{\partial f}{\partial s}(t,0)=J(t)$, we obtain that $J$ satisfies
the equation

$$
\frac{D^2 J}{dt^2}+R(\gamma'(t),J(t))\gamma'(t)=0.\qquad(1)
$$

The equation above is called the \emph{Jacobi equation}.

\begin{definition}[Jacobi field]
  \label{def:dc-ch5-2-1}
  \dcref{ch5:2.1}
  \lean{Riemannian.Jacobi.exists_jacobiField, Riemannian.Jacobi.jacobiField_eqOn}
  \leanok
  \uses{lem:dc-ch5-2-1-parallel-frame, lem:dc-ch5-2-1-ode, lem:dc-ch5-2-1-frame-reduction,
    lem:dc-ch5-2-1-frame-existence, lem:dc-ch5-2-1-real-curvature}
  Let $\gamma:[0,a]\to M$ be a geodesic in $M$. A vector field $J$ along $\gamma$ is
  said to be a \emph{Jacobi field} if it satisfies the Jacobi equation (1), for all
  $t\in[0,a]$.
  
  A Jacobi field is determined by its initial conditions $J(0)$,
  $\frac{DJ}{dt}(0)$. Indeed, let $e_1(t),\dots,e_n(t)$ be parallel, orthonormal
  fields along $\gamma$. Writing
  
  $$
  J(t)=\sum_i f_i(t)e_i(t),\quad a_{ij}=\langle R(\gamma'(t),e_i(t))\gamma'(t),e_j(t)\rangle,\quad i,j=1,\dots,n=\dim M,
  $$
  
  one has $\frac{D^2 J}{dt^2}=\sum_i f_i''(t)e_i(t)$ and
  $R(\gamma',J)\gamma'=\sum_{ij}f_i a_{ij}e_j$. Therefore equation (1) is equivalent
  to the system
  
  $$
  f_j''(t)+\sum_i a_{ij}(t)f_i(t)=0,\quad j=1,\dots,n,
  $$
  
  which is a linear system of the second order. Hence, given the initial
  conditions $J(0)$, $\frac{DJ}{dt}(0)$, there exists a $C^\infty$ solution defined
  on $[0,a]$. There exist, therefore, $2n$ linearly independent Jacobi fields along
  $\gamma$.
\end{definition}

The ingredients of the argument above are formalized in coordinate form:
the parallel orthonormal frame $e_1(t),\dots,e_n(t)$
(\cref{lem:dc-ch5-2-1-parallel-frame}); the existence/uniqueness of the
second-order linear system it reduces the Jacobi equation to
(\cref{lem:dc-ch5-2-1-ode}); the intrinsic reduction $J=\sum_i f_i e_i \Rightarrow
\frac{D^2J}{dt^2}=\sum_i f_i''e_i$ and the pairing with $e_j$
(\cref{lem:dc-ch5-2-1-frame-reduction,lem:dc-ch5-2-1-frame-existence}); and finally the
identification of the abstract coefficient with the genuine curvature contraction
$a_{ij}=\langle R(\gamma',e_i)\gamma',e_j\rangle$ read in the chart
(\cref{lem:dc-ch5-2-1-real-curvature}), which closes \cref{def:dc-ch5-2-1}.

\begin{lemma}[parallel orthonormal frame along a geodesic, coordinate form]
  \label{lem:dc-ch5-2-1-parallel-frame}
  \dcref{ch5:2.1}
  \lean{Riemannian.Jacobi.exists_parallelOrthoFrame_self,
    Riemannian.Jacobi.exists_parallelOrthoFrame,
    Riemannian.Jacobi.exists_chartOrthonormalBasis_self,
    Riemannian.Jacobi.chartMetricInner_const_of_parallelSol}
  \leanok
  \uses{def:dc-ch2-2-5, lem:dc-ch2-3-1-coord}
  Read in a coordinate chart at $\alpha$ along the coordinate curve $u$ of a geodesic
  $\gamma$ (with $\alpha=\gamma(0)$): the chart inner product $\langle\cdot,\cdot\rangle$ at
  the initial point is a positive-definite symmetric form on the model space, so it admits an
  orthonormal frame $e_1(0),\dots,e_n(0)$; parallel transport then produces coordinate vector
  fields $e_1(t),\dots,e_n(t)$ along $u$ that are parallel ($\frac{De_i}{dt}=0$) and remain
  orthonormal, $\langle e_i(t),e_j(t)\rangle=\delta_{ij}$, for every $t\in[0,a]$.
  Orthonormality is preserved because parallel transport is an isometry: parallel fields
  have constant inner product (the integrated form of \cref{lem:dc-ch2-3-1-coord}). These are
  do Carmo's ``parallel, orthonormal fields $e_1(t),\dots,e_n(t)$ along $\gamma$.''
\end{lemma}

\begin{lemma}[Jacobi system in a parallel frame: existence and uniqueness]
  \label{lem:dc-ch5-2-1-ode}
  \dcref{ch5:2.1}
  \lean{Riemannian.Jacobi.exists_isJacobiPairOn,
    Riemannian.Jacobi.IsJacobiPairOn.eqOn,
    Riemannian.Jacobi.IsJacobiPairOn.add,
    Riemannian.Jacobi.IsJacobiPairOn.smul}
  \leanok
  In a parallel orthonormal frame $e_1(t),\dots,e_n(t)$ along $\gamma$, writing
  $J(t)=\sum_i f_i(t)e_i(t)$ and $a_{ij}=\langle R(\gamma',e_i)\gamma',e_j\rangle$, the Jacobi
  equation (1) becomes the second-order linear system $f_j''(t)+\sum_i a_{ij}(t)f_i(t)=0$,
  i.e. $f''+A(t)f=0$ with $A(t)=(a_{ij}(t))$ the curvature contraction. For a continuous
  coefficient $A$ and any initial data $(f(0),f'(0))$ this system has a solution on $[0,a]$
  (do Carmo's ``there exists a $C^\infty$ solution''), determined uniquely by its initial
  conditions, and its solutions form a linear space closed under sum and scalar multiple (the
  $2n$ linearly independent Jacobi fields).
\end{lemma}

\begin{lemma}[Jacobi reduction in a parallel orthonormal frame]
  \label{lem:dc-ch5-2-1-frame-reduction}
  \dcref{ch5:2.1}
  \lean{Riemannian.Jacobi.covariantDerivCoord_frameCombination,
    Riemannian.Jacobi.covariantDerivCoord2_frameCombination,
    Riemannian.Jacobi.chartMetricInner_covariantDerivCoord2_frame,
    Riemannian.Jacobi.chartMetricInner_map_frameCombination_left,
    Riemannian.Jacobi.frameJacobiComponent,
    Riemannian.Jacobi.frameExpansion,
    Riemannian.Jacobi.covariantDerivCoord2_add_map_frameCombination_expand}
  \leanok
  \uses{lem:dc-ch5-2-1-parallel-frame, prop:dc-ch2-2-2}
  Write $J(t)=\sum_i f_i(t)e_i(t)$ in a parallel orthonormal frame
  $e_1(t),\dots,e_n(t)$ along $\gamma$ (\cref{lem:dc-ch5-2-1-parallel-frame}). Because
  the frame is parallel, $\frac{De_i}{dt}=0$, the covariant derivative differentiates
  only the scalar components (the Leibniz rule of \cref{prop:dc-ch2-2-2}):
  $$
  \frac{D}{dt}\sum_i f_i e_i=\sum_i f_i'\,e_i,\qquad
  \frac{D^2}{dt^2}\sum_i f_i e_i=\sum_i f_i''\,e_i.
  $$
  Pairing with $e_j$ and using orthonormality $\langle e_i,e_j\rangle=\delta_{ij}$
  extracts the $j$-th component, $\langle\frac{D^2J}{dt^2},e_j\rangle=f_j''$, while
  linearity of the curvature term in $J$ gives
  $\langle R(\gamma',J)\gamma',e_j\rangle=\sum_i f_i a_{ij}$ with
  $a_{ij}=\langle R(\gamma',e_i)\gamma',e_j\rangle$. Therefore pairing the Jacobi
  equation~(1) with $e_j$ is exactly do Carmo's scalar second-order system
  $f_j''+\sum_i a_{ij}f_i=0$; this is the reduction that turns \cref{def:dc-ch5-2-1}
  into the linear ODE of \cref{lem:dc-ch5-2-1-ode}.

  Conversely, the frame being orthonormal of full rank $n=\dim M$ is a basis, so every
  vector expands as $w=\sum_j\langle w,e_j\rangle e_j$, and the left-hand side of the
  Jacobi equation expands in the frame as
  $\frac{D^2J}{dt^2}+R(\gamma',J)\gamma'=\sum_j\big(f_j''+\sum_i a_{ij}f_i\big)e_j$.
  Thus~(1) holds if and only if the scalar system does: equation~(1) is \emph{equivalent}
  to the system, as do Carmo asserts.
\end{lemma}

\begin{lemma}[existence and uniqueness of a Jacobi field by its initial data, frame form]
  \label{lem:dc-ch5-2-1-frame-existence}
  \dcref{ch5:2.1}
  \lean{Riemannian.Jacobi.jacobiCoefOp,
    Riemannian.Jacobi.exists_jacobiField_frame,
    Riemannian.Jacobi.jacobiField_frame_eqOn}
  \leanok
  \uses{lem:dc-ch5-2-1-parallel-frame, lem:dc-ch5-2-1-ode, lem:dc-ch5-2-1-frame-reduction}
  Fix a parallel orthonormal frame $e_1(t),\dots,e_n(t)$ along $\gamma$
  (\cref{lem:dc-ch5-2-1-parallel-frame}), and let the curvature contraction
  $w\mapsto R(\gamma',w)\gamma'$ be read in the frame as the continuous coefficient field
  $a_{ij}(t)=\langle R(\gamma',e_i)\gamma',e_j\rangle$. Then, for every initial pair
  $\big(J(0),\tfrac{DJ}{dt}(0)\big)=(J_0,w_0)$, there exists a Jacobi field
  $J(t)=\sum_i f_i(t)e_i(t)$ along $\gamma$ with $J(0)=J_0$, initial velocity $w_0$, and
  $\frac{D^2J}{dt^2}+R(\gamma',J)\gamma'=0$; moreover any two Jacobi fields sharing this
  initial data coincide. This is do Carmo's assertion that ``given the initial conditions
  $J(0)$, $\frac{DJ}{dt}(0)$, there exists a $C^\infty$ solution defined on $[0,a]$,'' unique
  by its initial conditions. It reduces \cref{def:dc-ch5-2-1} to the identification of the
  contraction $w\mapsto R(\gamma',w)\gamma'$ with the intrinsic curvature operator read in the
  chart (the chart-curvature bridge), the last analytic input needed to close the definition
  itself.
\end{lemma}
\begin{proof}
  \leanok
  \uses{lem:dc-ch5-2-1-parallel-frame, lem:dc-ch5-2-1-ode, lem:dc-ch5-2-1-frame-reduction}
  Package the coefficients $a_{ij}(t)$ as a continuous operator field $A(t)$ acting on the
  coefficient vector $(f_i)$ by $(A(t)f)_j=\sum_i a_{ij}(t)f_i$. By \cref{lem:dc-ch5-2-1-ode}
  the second-order linear system $f''+A(t)f=0$ has a solution with initial data
  $f_i(0)=\langle J_0,e_i(0)\rangle$, $f_i'(0)=\langle w_0,e_i(0)\rangle$, unique given that
  data. Put $J=\sum_i f_i e_i$. Pairing the Jacobi operator with the frame
  (\cref{lem:dc-ch5-2-1-frame-reduction}) expands
  $\frac{D^2J}{dt^2}+R(\gamma',J)\gamma'=\sum_j\big(f_j''+\sum_i a_{ij}f_i\big)e_j=0$, so $J$
  is a Jacobi field; the orthonormal expansion gives
  $J(0)=\sum_i\langle J_0,e_i(0)\rangle e_i(0)=J_0$, and the frame being parallel gives the
  initial velocity $w_0$. Uniqueness of $J$ is uniqueness of the scalar solution read back
  through the frame. $\square$
\end{proof}

\begin{lemma}[the Jacobi coefficient is the curvature contraction read in the chart]
  \label{lem:dc-ch5-2-1-real-curvature}
  \dcref{ch5:2.1}
  \lean{Riemannian.Jacobi.chartCurvatureCoef,
    Riemannian.Jacobi.chartCurvatureCoef_contDiffOn,
    Riemannian.Jacobi.chartCurvatureOp,
    Riemannian.Jacobi.continuousOn_chartCurvatureOp,
    Riemannian.Jacobi.exists_jacobiField,
    Riemannian.Jacobi.jacobiField_eqOn}
  \leanok
  \uses{lem:dc-ch5-2-1-frame-existence, rem:dc-ch4-2-6}
  Read in the fixed chart at $\alpha=\gamma(0)$, the curvature contraction
  $w\mapsto R(\gamma',w)\gamma'$ along the coordinate curve $u$ (with velocity
  $\dot u=\gamma'$) is the operator field
  $$
  \big(R(t)\,w\big)^l=\sum_{i,j,k}R^l{}_{ijk}\big(u(t)\big)\,\dot u^i(t)\,w^j\,\dot u^k(t),
  \qquad
  R^l{}_{ijk}=\partial_j\Gamma^l_{ik}-\partial_i\Gamma^l_{jk}
    +\sum_s\big(\Gamma^s_{ik}\Gamma^l_{js}-\Gamma^s_{jk}\Gamma^l_{is}\big),
  $$
  where $R^l{}_{ijk}$ is the coordinate curvature coefficient of the fixed chart
  (\cref{rem:dc-ch4-2-6}), matching $R(\partial_i,\partial_j)\partial_k=\sum_l
  R^l{}_{ijk}\,\partial_l$. Because the Christoffel symbols are $C^\infty$, each
  $R^l{}_{ijk}$ is $C^\infty$ on the chart interior; hence along a $C^1$ geodesic curve
  $u$ whose image stays in the chart interior the operator field $t\mapsto R(t)$ is
  continuous. Feeding this genuine curvature contraction (in place of the abstract
  coefficient field) to \cref{lem:dc-ch5-2-1-frame-existence} yields, for every initial
  pair $(J_0,w_0)$, a Jacobi field $J=\sum_i f_i e_i$ with $J(a)=J_0$,
  $\tfrac{DJ}{dt}(a)=w_0$ and $\tfrac{D^2J}{dt^2}+R(\gamma',J)\gamma'=0$ on $(a,b)$,
  unique given its initial data. This is do Carmo's identification of the Jacobi
  coefficient $a_{ij}=\langle R(\gamma',e_i)\gamma',e_j\rangle$ with the intrinsic
  curvature operator read in the chart, closing \cref{def:dc-ch5-2-1}.
\end{lemma}

\begin{remark}[Remark 2.2]
  \label{rem:dc-ch5-2-2}
  \dcref{ch5:2.2}
  \lean{Riemannian.Jacobi.isJacobiFieldAlongOn_velocity,
    Riemannian.Jacobi.isJacobiFieldAlongOn_smul_velocity}
  \leanok
  \uses{def:dc-ch5-2-1}
  It is worth noting that $\gamma'(t)$ and $t\gamma'(t)$ are Jacobi fields along
  $\gamma$. The first field has derivative zero and vanishes nowhere; the second is
  zero if and only if $t=0$. Due to these facts, we shall consider Jacobi fields
  along $\gamma$ that are normal to $\gamma'$.

  Formalized (manifold form) with $\gamma'(t)=\mathrm{mfderiv}\,\gamma\,t\,1$ the own-foot
  velocity field, whose chart reading is the chart velocity $\dot u$. The pair
  $(\gamma', D\gamma'/dt)=(\dot u, 0)$ is a Jacobi pair because $\nabla\dot u=0$ (the
  fixed-chart geodesic ODE) and $\mathcal R(\dot u,\dot u)\dot u=0$ (curvature
  antisymmetry in the first pair, \lstinline{chartCurvature_self_fst}); the pair
  $(t\gamma', D(t\gamma')/dt)=(t\dot u, \dot u)$ likewise, since
  $\nabla(t\dot u)=\dot u+t\nabla\dot u=\dot u$ and
  $\mathcal R(t\dot u,\dot u)\dot u = t\,\mathcal R(\dot u,\dot u)\dot u=0$. The second
  field vanishes only at $t=0$; this is what forces $\gamma'(0)\notin\ker\Theta$ in
  \cref{rem:dc-ch5-3-2}.
\end{remark}

\begin{example}[Jacobi fields on manifolds of constant curvature]
  \label{ex:dc-ch5-2-3}
  \dcref{ch5:2.3}
  \lean{Riemannian.Jacobi.isJacobiFieldOn_of_constantCurvature,
    Riemannian.Jacobi.isJacobiFieldOn_constCurvatureSol_pos,
    Riemannian.Jacobi.isJacobiFieldOn_constCurvatureSol_zero,
    Riemannian.Jacobi.isJacobiFieldOn_constCurvatureSol_neg}
  \leanok
  \uses{lem:dc-ch4-3-4, lem:dc-ch5-2-3-const-form, lem:dc-ch5-2-3-jacobi-op}
  Let $M$ be a Riemannian manifold of constant sectional curvature $K$, and let
  $\gamma:[0,\ell]\to M$ be a normalized geodesic on $M$. Let $J$ be a Jacobi field
  along $\gamma$, normal to $\gamma'$. We claim that from $|\gamma'|=1$ and from
  \cref{lem:dc-ch4-3-4}, it follows that
  
  $$
  R(\gamma',J)\gamma'=KJ.
  $$
  
  Indeed, for all vector fields $T$ along $\gamma$,
  
  $$
  \langle R(\gamma',J)\gamma',T\rangle=K\{\langle\gamma',\gamma'\rangle\langle J,T\rangle-\langle\gamma',T\rangle\langle J,\gamma'\rangle\}=K\langle J,T\rangle.
  $$
  
  As a result, the Jacobi equation can be written as
  
  $$
  \frac{D^2 J}{dt^2}+KJ=0.\qquad(2)
  $$
  
  Let $w(t)$ be a parallel field along $\gamma$ with $\langle\gamma'(t),w(t)\rangle=0$
  and $|w(t)|=1$. It is easy to verify that
  
  $$
  J(t)=\begin{cases}\dfrac{\sin(t\sqrt{K})}{\sqrt{K}}\,w(t),&\text{if }K>0,\\[1ex] t\,w(t),&\text{if }K=0,\\[1ex]\dfrac{\sinh(t\sqrt{-K})}{\sqrt{-K}}\,w(t),&\text{if }K<0,\end{cases}
  $$
  
  is a solution of (2) with initial conditions $J(0)=0$, $J'(0)=w(0)$.
  
  As we saw previously, given $p\in M$, $v\in T_p M$, and $w\in T_v(T_p M)$, we can
  construct a Jacobi field along the geodesic $\gamma:[0,1]\to M$ given by
  $\gamma(t)=\exp_p tv$. For that, consider $f(t,s)=\exp_p tv(s)$, where $v(s)$ is a
  curve in $T_p M$ with $v(0)=v$, $v'(0)=w$, and take
  $J(t)=\frac{\partial f}{\partial s}(t,0)$. Notice that $J(0)=0$. We are going to
  show that this is essentially the only way of constructing Jacobi fields along
  $\gamma(t)$ with $J(0)=0$.

  \emph{Formalization.} The curvature reduction $R(\gamma',J)\gamma'=KJ$ (equation~(2)) is the
  content of \cref{lem:dc-ch5-2-3-const-form,lem:dc-ch5-2-3-jacobi-op}
  (\lstinline{chartCurvatureOp\_isConstantCurvature}). The claim that the three closed forms
  \emph{solve} (2) with $J(0)=0$, $J'(0)=w(0)$ is
  \lstinline{isJacobiFieldOn\_of\_constantCurvature}: for a parallel unit field $w$ normal to a
  unit-speed velocity (read in the fixed chart) and any scalar $h$ with $h''+Kh=0$, the pair
  $(h\,w,\,h'\,w)$ is a chart Jacobi field (\lstinline{IsJacobiFieldOn}) --- since $w$ is
  parallel, $D^2(hw)/dt^2=h''\,w$, and $R(\gamma',hw)\gamma'=hK\,w$ by the reduction, so (2)
  collapses to $(h''+Kh)\,w=0$. The three scalars $\sin(t\sqrt K)/\sqrt K$, $t$,
  $\sinh(t\sqrt{-K})/\sqrt{-K}$ solve $h''+Kh=0$ with $h(0)=0$, $h'(0)=1$
  (\lstinline{isJacobiFieldOn\_constCurvatureSol\_pos/zero/neg}), so $J(0)=0$ and $J'(0)=w(0)$.
\end{example}

\begin{lemma}[constant curvature: the pointwise $(0,4)$ form, coordinate form]
  \label{lem:dc-ch5-2-3-const-form}
  \dcref{ch5:2.3}
  \lean{Riemannian.Jacobi.curvatureFormAt_isConstantCurvature}
  \leanok
  \uses{lem:dc-ch4-3-4}
  Read at a point $q$, the field-level constant-curvature hypothesis
  $\langle R(X,Y)Z,W\rangle=K_0(\langle X,Z\rangle\langle Y,W\rangle-\langle Y,Z\rangle\langle X,W\rangle)$
  (\cref{lem:dc-ch4-3-4}, \lstinline{IsConstantCurvature}) gives the pointwise curvature $(0,4)$
  form $\langle R(x,y)z,t\rangle_g=K_0(\langle x,z\rangle\langle y,t\rangle-\langle y,z\rangle\langle x,t\rangle)$
  for all $x,y,z,t\in T_qM$.  Formalized by evaluating the hypothesis on smooth extensions of the
  four tangent vectors (\lstinline{curvatureOperatorAt_eq}).
\end{lemma}

\begin{lemma}[constant curvature: $R(\gamma',J)\gamma'=K J$, coordinate form]
  \label{lem:dc-ch5-2-3-jacobi-op}
  \dcref{ch5:2.3}
  \lean{Riemannian.Jacobi.chartCurvatureOp_isConstantCurvature}
  \leanok
  \uses{lem:dc-ch5-2-3-const-form, cor:dc-ch5-2-9}
  On a manifold of constant sectional curvature $K_0$, do Carmo's identity $R(\gamma',J)\gamma'=K_0J$
  for $J$ normal to a unit-speed geodesic.  Read in the fixed chart at $p$: for $w$ with
  $\langle w,\dot u\rangle=0$ and $\langle\dot u,\dot u\rangle=1$, the Jacobi-equation curvature
  operator satisfies $\mathcal R_{\mathrm{chart}}(\dot u,w)\dot u=K_0\,w$
  (\lstinline{chartCurvatureOp\_isConstantCurvature}).  Consequently the Jacobi equation reduces to
  the scalar system $D^2J/dt^2+K_0J=0$ (do Carmo's equation~(2)).  Proof: the operator is
  $\mathcal R_{\mathrm{chart}}(\dot u,w)\dot u=\text{chartCurvature}(\dot u)\,w\,\dot u\,\dot u$;
  pairing with an arbitrary $z$ through the manifold$\leftrightarrow$chart frame bridge
  (\lstinline{curvatureFormAt\_chartFrame}, as in \cref{cor:dc-ch5-2-9}) and substituting the
  constant-curvature form (\cref{lem:dc-ch5-2-3-const-form}), orthonormality collapses the model
  form to $K_0\langle w,z\rangle$; positive-definiteness of the chart Gram form
  (\lstinline{chartMetricInner\_pos}) then gives the operator identity.
\end{lemma}

\begin{proposition}[Proposition 2.4]
  \label{prop:dc-ch5-2-4}
  \dcref{ch5:2.4}
  \lean{Riemannian.Jacobi.expDifferential_eq_jacobiField,
    Riemannian.Jacobi.IsJacobiFieldAlongOn.eqOn_of_initial}
  \leanok
  \uses{def:dc-ch5-2-1}
  Let $\gamma:[0,a]\to M$ be a geodesic and let $J$ be a Jacobi field along
  $\gamma$ with $J(0)=0$. Put $\frac{DJ}{dt}(0)=w$ and $\gamma'(0)=v$. Using the
  natural identification $T_{av}(T_{\gamma(0)}M)\simeq T_{\gamma(0)}M$, construct a
  curve $v(s)$ in $T_{\gamma(0)}M$ with $v(0)=av$, $v'(0)=aw$. Put
  $f(t,s)=\exp_p(\frac{t}{a}v(s))$,
  $p=\gamma(0)$, and define a Jacobi field $\bar J$ by
  $\bar J(t)=\frac{\partial f}{\partial s}(t,0)$. Then $\bar J=J$ on $[0,a]$.
\end{proposition}
\begin{proof}
  \leanok
  \uses{def:dc-ch5-2-1}
  For $s=0$,
  
  $$
  \frac{D}{dt}\frac{\partial f}{\partial s}=\frac{D}{\partial t}\big((d\exp_p)_{tv}(tw)\big)
  =\frac{D}{\partial t}\big(t(d\exp_p)_{tv}(w)\big)
  =(d\exp_p)_{tv}(w)+t\frac{D}{\partial t}\big((d\exp_p)_{tv}(w)\big).
  $$
  
  Therefore, for $t=0$,
  
  $$
  \frac{D\bar J}{dt}(0)=\frac{D}{\partial t}\frac{\partial f}{\partial s}(0,0)=(d\exp_p)_o(w)=w.
  $$
  
  Since $J(0)=\bar J(0)=0$ and $\frac{DJ}{dt}(0)=\frac{D\bar J}{dt}(0)=w$, we
  conclude, from the uniqueness theorem, that $J=\bar J$. $\square$
\end{proof}

\begin{corollary}[Corollary 2.5]
  \label{cor:dc-ch5-2-5}
  \dcref{ch5:2.5}
  \lean{Riemannian.Jacobi.expDifferential_eq_jacobiField}
  \leanok
  \uses{prop:dc-ch5-2-4, def:dc-ch5-2-1}
  Let $\gamma:[0,a]\to M$ be a geodesic and put $p=\gamma(0)$. Then a Jacobi field
  $J$ along $\gamma$ with $J(0)=0$ is given by

  $$
  J(t)=(d\exp_p)_{t\gamma'(0)}(tJ'(0)),\quad t\in[0,a].
  $$

  Formalized on a complete Riemannian manifold in the endpoint form: for $v,Z\in T_pM$,
  reading $\exp_p$ in a chart $\zeta$ around $\exp_p v$, the differential $(d\exp_p)_v$ is
  defined and sends $Z$ to the chart reading of $Y_Z(1)$, where $Y_Z$ is \emph{the} Jacobi
  field along the geodesic $t\mapsto\exp_p(tv)$ with $Y_Z(0)=0$, $\frac{DY_Z}{dt}(0)=Z$
  (existence and uniqueness from \cref{def:dc-ch5-2-1}). Substituting $v=t\gamma'(0)$,
  $Z=J'(0)$ and rescaling the geodesic to unit parameter recovers do Carmo's
  $J(t)=(d\exp_p)_{t\gamma'(0)}(tJ'(0))$. The identification uses the variational transport
  of the geodesic flow (the differential of $\exp_p$ along the ray is the Jacobi variational
  field), which is do Carmo's \cref{prop:dc-ch5-2-4}.
\end{corollary}

\begin{remark}[Remark 2.6]
  \label{rem:dc-ch5-2-6}
  \dcref{ch5:2.6}
  \uses{prop:dc-ch5-2-4}
  It is possible to obtain a construction analogous to \cref{prop:dc-ch5-2-4} for Jacobi
  fields that do not satisfy the condition $J(0)=0$. Since we shall not use this
  fact, we leave its proof as an exercise (Exercise 2).
  
  From now on, for simplicity of notation, we put $\frac{DJ}{dt}=J'$,
  $\frac{D^2 J}{dt^2}=J''$, etc.
\end{remark}

\begin{lemma}[Taylor expansion of $|J|^2$ in a parallel frame (analytic core)]
  \label{lem:dc-ch5-2-7-taylor-ode}
  \dcref{ch5:2.7}
  \lean{Riemannian.Jacobi.norm_sq_jacobi_isLittleO,
    Riemannian.Jacobi.norm_sq_jacobi_isLittleO_local}
  \leanok
  \uses{def:dc-ch5-2-1}
  Read in a parallel orthonormal frame $e_1,\dots,e_n$ along $\gamma$, write $J=\sum_i f_i e_i$,
  so that $g(t)=|J(t)|^2=\langle f(t),f(t)\rangle$ is the ordinary squared Euclidean norm of the
  coefficient vector, and the Jacobi equation is the plain second-order linear ODE $f''=-A(t)f$
  with $A(t)=\big(\langle R(\gamma',e_i)\gamma',e_j\rangle\big)_{ij}$ the frame curvature (a
  $C^\infty$ coefficient). If $f(0)=0$, then with $w=f'(0)$
  $$
  g(t)=\langle w,w\rangle\,t^2-\tfrac13\langle w,A(0)w\rangle\,t^4+o(t^4),\qquad t\to 0.
  $$
  This is the analytic content of \cref{prop:dc-ch5-2-7}: the four Taylor coefficients
  $g(0)=g'(0)=g'''(0)=0$, $g''(0)=2\langle w,w\rangle$, $g''''(0)=-8\langle w,A(0)w\rangle$ are
  computed by the explicit derivative chain, and the remainder is Taylor's theorem with Peano
  (little-$o$) remainder. \emph{The parallel frame is what makes do Carmo's identity (4)
  automatic}: since the frame is parallel, covariant differentiation only hits the scalar
  coefficients, so $\frac{D}{dt}(R(\gamma',J)\gamma')(0)$ reduces to $R(\gamma',J')\gamma'(0)$ with
  no covariant derivative of the curvature tensor, precisely because the terms carrying a factor
  of $f(0)=0$ drop out. Substituting the frame identification
  $\langle w,A(0)w\rangle=\langle R(\gamma',w)\gamma',w\rangle=\langle R(v,w)v,w\rangle$ gives
  do Carmo's equation~(3). The coefficient is stated as $\langle w,w\rangle t^2$; with $|w|=1$
  (\cref{prop:dc-ch5-2-7}) this is $t^2$. Two Lean forms are provided: the global one
  (\lstinline{norm_sq_jacobi_isLittleO}, hypotheses on all of $\mathbb R$) and a \emph{local} one
  (\lstinline{norm_sq_jacobi_isLittleO_local}), where the ODE data $f,v,A$ is only
  $C^\infty$ and solves $f'=v$, $v'=-Af$ on an open interval $s\ni 0$. The little-$o$ lives at
  $\mathcal N(0)$, so the local form suffices for the manifold instantiation: the geodesic
  $t\mapsto\exp_p(tv)$ and its parallel orthonormal frame are $C^\infty$ only on the open time
  window where $\gamma$ stays in a chart around $p$.
\end{lemma}

\begin{lemma}[$C^\infty$ smoothness of the parallel frame and the Jacobi coefficient]
  \label{lem:dc-ch5-2-7-frame-smooth}
  \dcref{ch5:2.7}
  \lean{Riemannian.Jacobi.exists_contDiffOn_parallelOrthoFrame,
    Riemannian.Jacobi.contDiffOn_infty_jacobiCoefOp}
  \leanok
  \uses{lem:dc-ch5-2-1-parallel-frame, lem:dc-ch5-2-1-real-curvature}
  The regularity input for the manifold Taylor expansion (\cref{prop:dc-ch5-2-7}): along a curve
  $u$ that is $C^\infty$ on an open time interval and stays in the chart interior, the parallel
  orthonormal frame $e_1(t),\dots,e_n(t)$ can be chosen \emph{$C^\infty$} on the interior of a
  closed subinterval, and the Jacobi coefficient operator
  $A(t)=\big(a_{ij}(t)\big)=\big(\langle R(\gamma',e_i)\gamma',e_j\rangle\big)$ is then $C^\infty$
  in $t$.  The parallel-transport ODE $\dot e_i=-\Gamma(\dot u,e_i)(u)=B(t)e_i$ has a $C^\infty$
  operator coefficient once $u$ is $C^\infty$, so the linear-ODE $C^\infty$ bootstrap upgrades the
  frame from the $C^1$ regularity of \cref{lem:dc-ch5-2-1-parallel-frame} to $C^\infty$; the
  coefficient $a_{ij}(t)=\langle R(\dot u,e_i)\dot u,e_j\rangle$ is a chart Gram pairing of the
  $C^\infty$ frame, velocity, and curvature contraction (\cref{lem:dc-ch5-2-1-real-curvature}).
  An orthonormal frame at an arbitrary interior chart point lets the frame start off the pole, so
  that $t=0$ is an interior time and the two-sided expansion applies.
\end{lemma}

\begin{proposition}[Proposition 2.7]
  \label{prop:dc-ch5-2-7}
  \dcref{ch5:2.7}
  \lean{Riemannian.Jacobi.norm_sq_jacobi_frame_isLittleO}
  \leanok
  \uses{lem:dc-ch5-2-7-taylor-ode, lem:dc-ch5-2-7-frame-smooth, def:dc-ch5-2-1}
  % Formalized (frame/coordinate form, norm_sq_jacobi_frame_isLittleO): read the Jacobi field in
  % a C^∞ parallel orthonormal frame J = Σ f_i e_i, so |J(t)|^2_g = ‖f(t)‖^2 (orthonormality) and
  % the coefficients f solve the frame Jacobi system f'' = -A(t)f with A = jacobiCoefOp the genuine
  % frame curvature (C^∞ by lem:dc-ch5-2-7-frame-smooth). Transferring to EuclideanSpace feeds the
  % analytic core lem:dc-ch5-2-7-taylor-ode; ⟨V(0),A(0)V(0)⟩ = ⟨R(v,w)v,w⟩ via sum_jacobiCoef_quadratic.
  Let $p\in M$ and $\gamma:[0,a]\to M$ be a geodesic with $\gamma(0)=p$,
  $\gamma'(0)=v$. Let $w\in T_v(T_p M)$ with $|w|=1$ and let $J$ be a Jacobi field
  along $\gamma$ given by
  
  $$
  J(t)=(d\exp_p)_{tv}(tw),\quad 0\le t\le a.
  $$
  
  Then the Taylor expansion of $|J(t)|^2$ about $t=0$ is given by
  
  $$
  |J(t)|^2=t^2-\tfrac{1}{3}\langle R(v,w)v,w\rangle t^4+R(t),\qquad(3)
  $$
  
  where $\lim_{t\to 0}\frac{R(t)}{t^4}=0$.
\end{proposition}
\begin{proof}
  \leanok
  \uses{lem:dc-ch5-2-7-taylor-ode, lem:dc-ch5-2-7-frame-smooth}
  The Lean proof takes the frame route (\cref{lem:dc-ch5-2-7-taylor-ode}), where do Carmo's
  identity~(4) below is automatic: in a parallel orthonormal frame the covariant derivative hits
  only the scalar coefficients, and every term carrying a factor of $J(0)=0$ drops out, so no
  covariant derivative of the curvature tensor is needed.  The manifold reading
  $|J(t)|^2_g=\|f(t)\|^2$ and the coefficient identification
  $\langle J',R(\gamma',J')\gamma'\rangle(0)=\langle R(v,w)v,w\rangle$ are supplied by
  \cref{lem:dc-ch5-2-7-frame-smooth}.  do Carmo's original derivation follows.
  Since $J(0)=0$ and $J'(0)=w$, the first three coefficients are
  
  $$
  \langle J,J\rangle(0)=0,\quad\langle J,J\rangle'(0)=2\langle J,J'\rangle(0)=0,\quad\langle J,J\rangle''(0)=2\langle J',J'\rangle(0)+2\langle J'',J\rangle(0)=2.
  $$
  
  Since $J''(0)=-R(\gamma',J)\gamma'(0)=0$, we have
  $\langle J,J\rangle'''(0)=6\langle J',J''\rangle(0)+2\langle J''',J\rangle(0)=0$. Now
  we need the fact
  
  $$
  \nabla_{\gamma'}(R(\gamma',J)\gamma')(0)=R(\gamma',J')\gamma'(0).\qquad(4)
  $$
  
  To prove (4), note that for any $W$, at $t=0$,
  
  $$
  \Big\langle\frac{D}{dt}(R(\gamma',J)\gamma'),W\Big\rangle
  =\frac{d}{dt}\langle R(\gamma',W)\gamma',J\rangle-\langle R(\gamma',J)\gamma',W'\rangle
  =\Big\langle\frac{D}{dt}(R(\gamma',W)\gamma'),J\Big\rangle+\langle R(\gamma',W)\gamma',J'\rangle
  =\langle R(\gamma',J')\gamma',W\rangle,
  $$
  
  which implies (4). It follows from (4) and the Jacobi equation that
  $J'''(0)=-R(\gamma',J')\gamma'(0)$. Therefore,
  
  $$
  \langle J,J\rangle''''(0)=8\langle J',J'''\rangle(0)+6\langle J'',J''\rangle(0)+2\langle J'''',J\rangle(0)=-8\langle J',R(\gamma',J')\gamma'\rangle(0)=-8\langle R(v,w)v,w\rangle.
  $$
  
  Putting together the calculation above, we obtain (3). $\square$
\end{proof}

\begin{remark}[Remark 2.8]
  \label{rem:dc-ch5-2-8}
  \dcref{ch5:2.8}
  The expression (4) can also be obtained using the (covariant) derivation of
  tensors described in Section 5 of Chap. 4 (see Exercise 5 of Chap. 4).
\end{remark}

\begin{corollary}[Corollary 2.9]
  \label{cor:dc-ch5-2-9}
  \dcref{ch5:2.9}
  \lean{Riemannian.Jacobi.chartMetricInner_chartCurvatureOp_eq_sectionalCurvature,
    Riemannian.Jacobi.chartMetricInner_chartCurvatureOp_eq_curvatureFormAt}
  \leanok
  \uses{def:dc-ch4-3-2, prop:dc-ch5-2-7}
  If $\gamma:[0,\ell]\to M$ is parametrized by arc length (i.e. $|v|=1$) and
  $\langle w,v\rangle=0$, the expression $\langle R(v,w)v,w\rangle$ is the sectional
  curvature at $p$ with respect to the plane $\sigma$ generated by $v$ and $w$.
  Therefore, in this situation,

  $$
  |J(t)|^2=t^2-\tfrac{1}{3}K(p,\sigma)t^4+R(t).\qquad(5)
  $$

  Formalized (coordinate form, \lstinline{chartMetricInner_chartCurvatureOp_eq_sectionalCurvature}):
  the Jacobi coefficient $\langle R(v,w)v,w\rangle$ carried by \cref{prop:dc-ch5-2-7} is read in
  the fixed chart at $p$ as
  $\langle\mathcal R_{\mathrm{chart}}(\dot u,w)\dot u,w\rangle_G$ with $\dot u$ the chart velocity
  ($=\gamma'$).  The chart curvature operator is the manifold curvature $(0,4)$-form of the frame
  realizations $F v,F w\in T_qM$ ($q$ the manifold point, $\varphi(q)=u(t_0)$), via the
  operator$\leftrightarrow$vector reindexing (\lstinline{chartCurvatureOp_eq_chartCurvature}), the
  manifold$\leftrightarrow$chart bridge (\lstinline{curvatureFormAt_chartFrame}) and antisymmetry
  in the first pair (\lstinline{chartMetricInner_chartCurvatureOp_eq_curvatureFormAt}):
  $\langle R(v,w)v,w\rangle = \mathrm{curvatureForm}_q(Fv,Fw,Fv,Fw)$.  Dividing by the wedge
  denominator $\langle Fv,Fv\rangle\langle Fw,Fw\rangle-\langle Fv,Fw\rangle^2$, which the
  orthonormality hypotheses $|v|=|w|=1$, $\langle v,w\rangle=0$ (in chart Gram form, transferred to
  the intrinsic inner product of the realizations by \lstinline{metricInner_chartFrameRealize})
  force to $1$, gives exactly the sectional curvature
  $K(p,\sigma)=\mathrm{sectionalCurvature}(Fv,Fw)$ of \cref{def:dc-ch4-3-2}.  Feeding this into
  \cref{lem:dc-ch5-2-10-sqrt} produces do Carmo's $|J(t)|=t-\tfrac16 K(p,\sigma)t^3+o(t^3)$
  (\cref{cor:dc-ch5-2-10}).
\end{corollary}

\begin{lemma}[square root of the $|J|^2$ expansion (analytic core)]
  \label{lem:dc-ch5-2-10-sqrt}
  \dcref{ch5:2.10}
  \lean{Riemannian.Jacobi.sqrt_isLittleO_of_sq_isLittleO}
  \leanok
  \uses{cor:dc-ch5-2-9}
  Let $g\ge 0$ (playing the role of $|J(t)|^2$) have the fourth-order expansion
  $g(t)=t^2-c\,t^4+o(t^4)$ as $t\to 0$ (do Carmo's (5), with $c=\tfrac13 K(p,\sigma)$ after the
  sectional-curvature identification of \cref{cor:dc-ch5-2-9}). Then its square root has the
  third-order expansion
  $$
  \sqrt{g(t)}=t-\tfrac{c}{2}\,t^3+o(t^3),\qquad t\to 0^+.
  $$
  With $c=\tfrac13 K$ this is do Carmo's $|J(t)|=t-\tfrac16 K(p,\sigma)t^3+\tilde R(t)$ (his (6)),
  since $|J(t)|=\sqrt{|J(t)|^2}$. The odd-power expansion is one-sided ($t\ge 0$): $\sqrt{t^2}=|t|$
  agrees with $t$ only for $t\ge 0$, which is do Carmo's range $t\in[0,\ell]$. Proof: writing
  $p(t)=t-\tfrac{c}{2}t^3$, the numerator $g-p^2=(g-(t^2-ct^4))-\tfrac{c^2}{4}t^6=o(t^4)$, and
  $\sqrt g-p=(g-p^2)/(\sqrt g+p)$ with denominator $\ge t/2>0$ near $0^+$, so the quotient is
  $o(t^3)$.
\end{lemma}

\begin{corollary}[Corollary 2.10]
  \label{cor:dc-ch5-2-10}
  \dcref{ch5:2.10}
  \lean{Riemannian.Jacobi.norm_jacobi_frame_isLittleO}
  \leanok
  \uses{prop:dc-ch5-2-7, lem:dc-ch5-2-10-sqrt, cor:dc-ch5-2-9}
  With the same conditions as in the previous corollary,

  $$
  |J(t)|=t-\tfrac{1}{6}K(p,\sigma)t^3+\tilde R(t),\quad\lim_{t\to 0}\frac{\tilde R(t)}{t^3}=0.\qquad(6)
  $$

  Formalized (frame/coordinate form, \lstinline{norm_jacobi_frame_isLittleO}) with the
  normalization $|w|_g=1$: taking the square root of \cref{prop:dc-ch5-2-7} via the analytic
  $\sqrt{}$ step \cref{lem:dc-ch5-2-10-sqrt} gives the one-sided ($t\to 0^+$) expansion
  $|J(t)|_g = t - \tfrac16\langle R(v,w)v,w\rangle t^3 + o(t^3)$; the coefficient
  $\langle R(v,w)v,w\rangle$ is the sectional curvature $K(p,\sigma)$ by \cref{cor:dc-ch5-2-9}
  when $|v|=1$, $\langle v,w\rangle=0$.  The global nonnegativity the $\sqrt{}$ step needs is
  arranged by clamping $|J|^2_g$ (nonnegative on the chart) with $\max(\cdot,0)$.
\end{corollary}

\begin{remark}[Remark 2.11]
  \label{rem:dc-ch5-2-11}
  \dcref{ch5:2.11}
  The expression (6) essentially contains the relation between geodesics and
  curvature mentioned at the beginning of this chapter. Considering the
  parametrized surface
  
  $$
  f(t,s)=\exp_p tv(s),\quad t\in[0,\delta],\quad s\in(-\varepsilon,\varepsilon),
  $$
  
  where $\delta$ is chosen so small that $\exp_p tv(s)$ is defined, and $v(s)$ is a
  curve in $T_p M$ with $|v(s)|=1$, $v(0)=v$, $v'(0)=w$, we see that the rays
  $t\mapsto tv(s)$, $t\in[0,\delta]$, that start from the origin $0$ of $T_p M$,
  deviate from the ray $t\mapsto tv(0)$ with velocity
  
  $$
  \Big|\Big(\frac{\partial}{\partial s}tv(s)\Big)(0)\Big|=|tw|=t.
  $$
  
  On the other hand, (6) tells us that the geodesics $t\mapsto\exp_p(tv(s))$
  deviate from the geodesic $\gamma(t)=\exp_p tv(0)$ with a velocity that differs
  from $t$ by a term of third order in $t$, given by $-\tfrac{1}{6}K(p,\sigma)t^3$.
  This tells us that, locally, the geodesics spread apart less than the rays in
  $T_p M$ if $K_p(\sigma)>0$, and more than the rays in $T_p M$ if $K_p(\sigma)<0$.
  For $t$ small, the value $K(p,\sigma)t^3$ furnishes an approximation for the
  extent of this spread with an error of order $t^3$.
\end{remark}

\section{Conjugate points}

We now turn to the relationship between the singularities of the exponential map
and Jacobi fields.

\begin{definition}[conjugate point, multiplicity]
  \label{def:dc-ch5-3-1}
  \dcref{ch5:3.1}
  \lean{Riemannian.Jacobi.IsConjugatePointAt}
  \leanok
  \uses{def:dc-ch5-2-1}
  Let $\gamma:[0,a]\to M$ be a geodesic. The point $\gamma(t_0)$ is said to be
  \emph{conjugate} to $\gamma(0)$ along $\gamma$, $t_0\in(0,a]$, if there exists a Jacobi
  field $J$ along $\gamma$, not identically zero, with $J(0)=0=J(t_0)$. The maximum
  number of such linearly independent fields is called the \emph{multiplicity} of the
  conjugate point $\gamma(t_0)$.

  The predicate \lstinline{IsConjugatePointAt} formalizes the conjugacy relation exactly:
  there is a manifold Jacobi field $(J,DJ)$ along $\gamma|_{[0,t_0]}$ with $J(0)=J(t_0)=0$
  that is not identically zero. The multiplicity is realized as the dimension of
  $\ker(d\exp_p)_{t_0\gamma'(0)}$ in \cref{prop:dc-ch5-3-5}. The symmetry observation
  (if $\gamma(t_0)$ is conjugate to $\gamma(0)$ then $\gamma(0)$ is conjugate to
  $\gamma(t_0)$) is immediate from the symmetric roles of the two endpoints in the
  vanishing conditions.
\end{definition}

\begin{remark}[Remark 3.2]
  \label{rem:dc-ch5-3-2}
  \dcref{ch5:3.2}
  \lean{Riemannian.Jacobi.conjugateMultiplicity,
    Riemannian.Jacobi.conjugateMultiplicity_le_finrank,
    Riemannian.Jacobi.conjugateMultiplicity_le_finrank_sub_one}
  \leanok
  \uses{rem:dc-ch5-2-2}
  If $\dim M=n$, there exist exactly $n$ linearly independent Jacobi fields along
  the geodesic $\gamma:[0,a]\to M$ which are zero at $\gamma(0)$. This follows from
  the fact, easily checked, that the Jacobi fields $J_1,\dots,J_k$ with $J_i(0)=0$
  are linearly independent if and only if $J_1'(0),\dots,J_k'(0)$ are linearly
  independent. In addition, the Jacobi field $J(t)=t\gamma'(t)$ never vanishes for
  $t\neq 0$ (see \cref{rem:dc-ch5-2-2}). From this we deduce that the multiplicity of a
  conjugate point never exceeds $n-1$.

  Formalized (coordinate/endpoint form) via the initial-velocity parametrization of
  \cref{prop:dc-ch5-3-9}: Jacobi fields with $J(0)=0$ correspond bijectively to their
  initial velocity $J'(0)\in E$, so their space is $n$-dimensional, and the
  multiplicity $\dim\ker\Theta$ (\lstinline{conjugateMultiplicity}) is at most $n$
  (\lstinline{conjugateMultiplicity_le_finrank}). The sharp bound
  ``$\le n-1$'' (\lstinline{conjugateMultiplicity_le_finrank_sub_one}) uses that
  $t\mapsto t\gamma'(t)$ is a Jacobi field with $J(0)=0$ and initial velocity
  $\gamma'(0)$, whose value at the endpoint is $\Theta(\gamma'(0))=L\gamma'(L)\neq0$
  for a non-constant geodesic ($\gamma'(L)\neq0$; \cref{rem:dc-ch5-2-2}), so
  $\gamma'(0)\notin\ker\Theta$, hence $\ker\Theta\subsetneq E$ and
  $\dim\ker\Theta\le n-1$.
\end{remark}

\begin{example}[Example 3.3]
  \label{ex:dc-ch5-3-3}
  \dcref{ch5:3.3}
  \lean{Riemannian.Jacobi.isConjugatePointAt_constCurvature_pos,
    Riemannian.Jacobi.conjugateMultiplicity_constCurvature_pos_eq}
  \leanok
  \uses{ex:dc-ch5-2-3, def:dc-ch5-3-1, rem:dc-ch5-3-2, prop:dc-ch5-3-5}
  Let $S^n=\{x\in\mathbb{R}^{n+1};|x|=1\}$. Assuming the fact (to be proved in the
  next chapter) that the sectional curvatures of $S^n$ are all equal to one, the
  Jacobi field on $S^n$ given in \cref{ex:dc-ch5-2-3}, namely $J(t)=(\sin t)w(t)$, satisfies
  $J(0)=J(\pi)=0$. Therefore, along any geodesic $\gamma$ of $S^n$, the antipodal
  point $\gamma(\pi)$ to $\gamma(0)$ is conjugate to $\gamma(0)$. There exist $n-1$
  such fields that are linearly independent, that is, the multiplicity of
  $\gamma(\pi)$ as a conjugate point of $\gamma(0)$ is $n-1$.

  This is formalized as the intrinsic statement of which $S^n$ (the $K_0=1$ case) is
  an instance: on any manifold $M^n$ of constant sectional curvature $K_0>0$, along a
  unit-speed geodesic $\gamma\colon[0,\pi/\sqrt{K_0}]\to M$, the point
  $\gamma(\pi/\sqrt{K_0})$ is conjugate to $\gamma(0)$, with multiplicity $n-1$. The
  normal Jacobi fields with $J(0)=0$ are the explicit fields
  $J(t)=\frac{\sin(t\sqrt{K_0})}{\sqrt{K_0}}\,w(t)$ of \cref{ex:dc-ch5-2-3}, where $w$
  is the parallel transport along $\gamma$ of a vector $w_0\perp\gamma'(0)$. Because
  parallel transport is an isometry and $\gamma'$ is itself parallel,
  $\langle w(t),\gamma'(t)\rangle=\langle w_0,\gamma'(0)\rangle=0$ throughout, so $J$ is
  a Jacobi field with $J(0)=0$, $J'(0)=w_0$, and $J(\pi/\sqrt{K_0})=0$ since
  $\sin\pi=0$. By uniqueness of Jacobi fields with given initial data, every
  $w_0\perp\gamma'(0)$ lies in the kernel of the endpoint map
  $\Theta\colon J'(0)\mapsto J(\pi/\sqrt{K_0})$ (\cref{prop:dc-ch5-3-5}); hence
  $\gamma'(0)^\perp\subseteq\ker\Theta$ and $n-1=\dim\gamma'(0)^\perp\le\dim\ker\Theta$,
  while the reverse inequality is the sharp bound $\dim\ker\Theta\le n-1$ of
  \cref{rem:dc-ch5-3-2}, giving equality. The concrete $S^n$ statement is the $K_0=1$
  instance, using additionally the Ch.~6 fact that $S^n$ has constant sectional
  curvature $1$.

  The manifold-level parallel field $w(t)$ — parallel transport of a single vector
  along a geodesic that may leave any single chart, the situation of the antipodal arc
  from $p$ to $-p$ on $S^n$ — is built in \lstinline{ParallelFieldAlong.lean}
  (\lstinline{exists_parallelFieldAlongOn}), by a supremum walk gluing single-chart
  parallel fields through the chart-change covariance of the parallel-transport ODE;
  parallel transport is an isometry (\lstinline{IsParallelFieldAlongOn.metricInner_const}),
  and the manifold constant-curvature Jacobi field
  $\frac{\sin(t\sqrt{K_0})}{\sqrt{K_0}}w$ is
  \lstinline{isJacobiFieldAlongOn_constCurvatureSol_pos}.
\end{example}

\begin{definition}[conjugate locus]
  \label{def:dc-ch5-3-4}
  \dcref{ch5:3.4}
  The set of (first) conjugate points to the point $p\in M$, for all the geodesics
  that start at $p$, is called the \emph{conjugate locus} of $p$ and is denoted by
  $C(p)$.
  
  On $S^n$, $C(p)=\{-p\}$, for all $p$. The case of $S^n$, however, is not typical.
  A more typical example is given by the ellipsoid, where $C(p)$ is, in general, a
  curve with four singular points (see Fig. 4 of Chap. 13; cf. Braunmühl, A.,
  "Geodätische Linien auf dreiachsigen Flächen 2-Grades", Math. Ann. 20 (1882),
  557-586).
\end{definition}

\begin{proposition}[Proposition 3.5]
  \label{prop:dc-ch5-3-5}
  \dcref{ch5:3.5}
  \lean{Riemannian.Jacobi.expDifferential_injective_iff_not_conjugate,
    Riemannian.Jacobi.injective_jacobiEndpointOfVel_iff_not_conjugate,
    Riemannian.Jacobi.conjugateMultiplicity,
    Riemannian.Jacobi.isConjugatePointAt_iff_conjugateMultiplicity_pos,
    Riemannian.Jacobi.conjugateMultiplicity_eq_finrank_ker_expDifferential}
  \leanok
  \uses{cor:dc-ch5-2-5, def:dc-ch5-3-1}
  Let $\gamma:[0,a]\to M$ be a geodesic and put $\gamma(0)=p$. The point
  $q=\gamma(t_0)$, $t_0\in(0,a]$, is conjugate to $p$ along $\gamma$ if and only if
  $v_0=t_0\gamma'(0)$ is a critical point of $\exp_p$. In addition, the multiplicity
  of $q$ as a conjugate point of $p$ is equal to the dimension of the kernel of the
  linear map $(d\exp_p)_{v_0}$.

  The first assertion is formalized (axiom-clean) in the endpoint form $t_0=1$: on a
  complete manifold $(d\exp_p)_v$ is injective if and only if $\exp_p(v)=\gamma_v(1)$ is
  \emph{not} conjugate to $p$ along $\gamma_v$ (\cref{def:dc-ch5-3-1}). Equivalently,
  $v$ is a critical point of $\exp_p$ (its differential fails to be injective) exactly when
  the endpoint is conjugate. Both directions use \cref{cor:dc-ch5-2-5}: a nonzero
  $Z\in\ker(d\exp_p)_v$ yields the nonzero Jacobi field $Y_Z$ with $Y_Z(0)=Y_Z(1)=0$, and
  conversely. The equivalence is also available in endpoint-general form
  (\lstinline{injective_jacobiEndpointOfVel_iff_not_conjugate}): the endpoint map
  $\Theta:J'(0)\mapsto J(a)$ on Jacobi fields with $J(0)=0$ is injective iff
  $\gamma(a)$ is not conjugate to $\gamma(0)$. The multiplicity is \emph{defined}
  (\lstinline{conjugateMultiplicity}) as $\dim\ker\Theta$, the dimension of the
  space of Jacobi fields vanishing at both endpoints, and $\gamma(a)$ is conjugate
  iff this multiplicity is positive
  (\lstinline{isConjugatePointAt_iff_conjugateMultiplicity_pos}). The second
  assertion is also formalized (endpoint form,
  \lstinline{conjugateMultiplicity_eq_finrank_ker_expDifferential}): writing the
  differential $D=(d\exp_p)_v$ in a chart, $DZ$ is the chart reading of $J_{(0,Z)}(1)$,
  so $DZ=0$ iff $\Theta Z=0$ (the tangent chart reading is injective). Hence
  $\ker D=\ker\Theta$ and the multiplicity $\dim\ker\Theta$ equals $\dim\ker(d\exp_p)_{v}$.
\end{proposition}
\begin{proof}
  The point $q=\gamma(t_0)$ is a conjugate point of $p$ along $\gamma$ if
  and only if there exists a non-zero Jacobi field $J$ along $\gamma$ with
  $J(0)=J(t_0)=0$. Let $v=\gamma'(0)$ and $w=J'(0)$. From \cref{cor:dc-ch5-2-5},
  $J(t)=(d\exp_p)_{tv}(tw)$, $t\in[0,a]$. Observe that $J$ is non-zero if and only
  if $w\neq 0$. Therefore, $q=\gamma(t_0)$ is conjugate to $p$ if and only if
  
  $$
  0=J(t_0)=(d\exp_p)_{t_0 v}(t_0 w),\quad w\neq 0,
  $$
  
  that is, if and only if $t_0 v$ is a critical point of $\exp_p$. The first
  assertion is therefore proved.
  
  The multiplicity of $q$ is equal to the number of linearly independent Jacobi
  fields $J_1,\dots,J_k$ which are zero at $0$ and at $t_0$. The fields
  $J_1,\dots,J_k$ are linearly independent if and only if
  $J_1'(0),\dots,J_k'(0)$ are linearly independent in $T_p M$. From the
  construction above, the multiplicity of $q$ is equal to the dimension of the
  kernel of $(d\exp_p)_{t_0 v}$. $\square$
\end{proof}

\begin{proposition}[Proposition 3.6]
  \label{prop:dc-ch5-3-6}
  \dcref{ch5:3.6}
  \lean{Riemannian.Jacobi.chartMetricInner_jacobi_velocity_affine}
  \leanok
  \uses{def:dc-ch5-2-1}
  Let $J$ be a Jacobi field along the geodesic $\gamma:[0,a]\to M$. Then

  $$
  \langle J(t),\gamma'(t)\rangle=\langle J'(0),\gamma'(0)\rangle t+\langle J(0),\gamma'(0)\rangle,\quad t\in[0,a].
  $$

  Formalized in coordinate form: reading $\gamma$ in a chart $\alpha$ whose source contains
  $\gamma([a,b])$, with chart curve $u=\varphi_\alpha\circ\gamma$ and chart velocity
  $\dot u=\frac{d}{dt}u$ (the chart reading of $\gamma'$), the chart Gram pairing
  $\langle J,\dot u\rangle_G$ is an affine function of $t$:
  $\langle J(t),\dot u(t)\rangle=\langle DJ(a),\dot u(a)\rangle(t-a)+\langle J(a),\dot u(a)\rangle$.
  Since the chart Gram pairing equals the intrinsic $\langle\cdot,\cdot\rangle_g$ and $\dot u$ is
  the chart reading of $\gamma'$, this is do Carmo's identity. The proof follows do Carmo:
  $\langle J,\dot u\rangle'=\langle DJ,\dot u\rangle$ (metric compatibility, $\nabla J=DJ$, and the
  velocity is parallel $\nabla\dot u=0$ in the fixed chart — the fixed-chart geodesic ODE), and
  $\langle DJ,\dot u\rangle'=-\langle R(J,\dot u)\dot u,\dot u\rangle=0$ (curvature antisymmetry in
  the last pair), so $\langle DJ,\dot u\rangle$ is constant and $\langle J,\dot u\rangle$ is affine.
\end{proposition}
\begin{proof}
  \leanok
  \uses{def:dc-ch5-2-1}
  Omitting the $t$ for notation, we have from the Jacobi equation,
  
  $$
  \langle J',\gamma'\rangle'=\langle J'',\gamma'\rangle=-\langle R(\gamma',J)\gamma',\gamma'\rangle=0.
  $$
  
  Therefore, $\langle J',\gamma'\rangle=\langle J'(0),\gamma'(0)\rangle$. In
  addition,
  
  $$
  \langle J,\gamma'\rangle'=\langle J',\gamma'\rangle=\langle J'(0),\gamma'(0)\rangle.
  $$
  
  Integrating this last equation in $t$, we obtain finally
  $\langle J,\gamma'\rangle=\langle J'(0),\gamma'(0)\rangle t+\langle J(0),\gamma'(0)\rangle$.
  $\square$
\end{proof}

\begin{corollary}[Corollary 3.7]
  \label{cor:dc-ch5-3-7}
  \dcref{ch5:3.7}
  \lean{Riemannian.Jacobi.chartMetricInner_jacobi_velocity_const_of_eq}
  \leanok
  \uses{prop:dc-ch5-3-6}
  If $\langle J,\gamma'\rangle(t_1)=\langle J,\gamma'\rangle(t_2)$, $t_1,t_2\in[0,a]$,
  $t_1\neq t_2$, then $\langle J,\gamma'\rangle$ does not depend on $t$; in
  particular, if $J(0)=J(a)=0$, then $\langle J,\gamma'\rangle(t)\equiv 0$.

  Coordinate form (from the affine law of \cref{prop:dc-ch5-3-6}): if the chart Gram pairing
  $\langle J,\dot u\rangle$ agrees at two distinct times, its slope $\langle DJ(a),\dot u(a)\rangle$
  vanishes, so the pairing is constant. The ``in particular'' clause is immediate: $J(a)=0$ forces
  $\langle J(a),\dot u(a)\rangle=0$, hence the constant value is $0$.
\end{corollary}

\begin{corollary}[Corollary 3.8]
  \label{cor:dc-ch5-3-8}
  \dcref{ch5:3.8}
  \lean{Riemannian.Jacobi.chartMetricInner_jacobi_velocity_eq_zero_iff,
    Riemannian.Jacobi.velocityFunctional,
    Riemannian.Jacobi.finrank_velocityPerp_eq}
  \leanok
  \uses{prop:dc-ch5-3-6, rem:dc-ch5-3-2}
  Suppose that $J(0)=0$. Then $\langle J'(0),\gamma'(0)\rangle=0$ if and only if
  $\langle J,\gamma'\rangle(t)\equiv 0$; in particular, the space of Jacobi fields
  $J$ with $J(0)=0$ and $\langle J,\gamma'\rangle(t)\equiv 0$ has dimension equal to
  $n-1$.

  The equivalence is formalized (coordinate form, from \cref{prop:dc-ch5-3-6}): with $J(a)=0$
  the affine law reads $\langle J(t),\dot u(t)\rangle=\langle DJ(a),\dot u(a)\rangle(t-a)$, which
  vanishes identically iff its slope $\langle DJ(a),\dot u(a)\rangle=\langle J'(0),\gamma'(0)\rangle$
  does. The ``dimension $=n-1$'' clause is formalized as
  \lstinline{finrank_velocityPerp_eq}: under the initial-velocity parametrization
  ($w=J'(0)$) the space $\mathcal J^\perp$ of such $J$ is the kernel of the tangential
  functional $w\mapsto\langle w,\gamma'(0)\rangle_g$ (\lstinline{velocityFunctional}),
  which is a nonzero element of the dual for $\gamma'(0)\neq0$, hence a hyperplane of
  dimension $n-1$.
\end{corollary}

\begin{proposition}[Proposition 3.9]
  \label{prop:dc-ch5-3-9}
  \dcref{ch5:3.9}
  \lean{Riemannian.Jacobi.exists_unique_jacobi_of_endpoints}
  \leanok
  \uses{cor:dc-ch5-2-5, def:dc-ch5-3-1}
  Let $\gamma:[0,a]\to M$ be a geodesic. Let $V_1\in T_{\gamma(0)}M$ and
  $V_2\in T_{\gamma(a)}M$. If $\gamma(a)$ is not conjugate to $\gamma(0)$ there
  exists a unique Jacobi field $J$ along $\gamma$, with $J(0)=V_1$ and $J(a)=V_2$.

  Formalized in the endpoint form (base $0$, terminal $a$). Jacobi fields along
  $\gamma$ are parametrized by their initial data $(J(0),J'(0))\in E\times E$
  (existence and uniqueness from \cref{cor:dc-ch5-2-5} make this a bijection), so
  evaluation at a fixed time is a linear map $\mathrm{evalJacobiAt}:E\times E\to E$,
  and $\Theta=\mathrm{jacobiEndpointOfVel}:J'(0)\mapsto J(a)$ is its restriction to
  the fields with $J(0)=0$. By \cref{prop:dc-ch5-3-5} (endpoint form) $\Theta$ is
  injective iff $\gamma(a)$ is not conjugate to $\gamma(0)$; as $\dim M<\infty$,
  injectivity upgrades to surjectivity, and solving $\Theta w=V_2-\mathrm{evalJacobiAt}(V_1,0)$
  produces the field with $J(0)=V_1$, $J(a)=V_2$. Uniqueness: a competitor differs by
  a Jacobi field with $J(0)=J(a)=0$, i.e.\ an element of $\ker\Theta=0$.
\end{proposition}
\begin{proof}
  \leanok
  \uses{cor:dc-ch5-2-5, def:dc-ch5-3-1}
  Let $\mathcal{J}$ be the space of Jacobi fields $J$ with $J(0)=0$. Define
  a mapping $\Theta:\mathcal{J}\to T_{\gamma(a)}M$ by $\Theta(J)=J(a)$,
  $J\in\mathcal{J}$. Since $\gamma(a)$ is not conjugate to $\gamma(0)$, $\Theta$ is
  injective. Indeed, if $J_1\neq J_2$ with $J_1(a)=J_2(a)$, we should have
  $J_1-J_2$, a non-zero Jacobi field with $(J_1-J_2)(0)=0$, which contradicts the
  fact that $\gamma(a)$ is not conjugate to $\gamma(0)$.
  
  Since $\Theta$ is linear, it follows from injectivity and the fact that
  $\dim\mathcal{J}=\dim T_{\gamma(a)}M$ that $\Theta$ is an isomorphism. Hence there
  exists $\bar J_1\in\mathcal{J}$ with $\bar J_1(0)=0$ and $\bar J_1(a)=V_2$. By an
  analogous argument, there exists a Jacobi field $\bar J_2$ along $\gamma$ with
  $\bar J_2(a)=0$, $\bar J_2(0)=V_1$. The desired field is now given by
  $J=\bar J_1+\bar J_2$. Uniqueness is clear. $\square$
\end{proof}

\begin{corollary}[Corollary 3.10]
  \label{cor:dc-ch5-3-10}
  \dcref{ch5:3.10}
  \lean{Riemannian.Jacobi.jacobiConjugateBasis,
    Riemannian.Jacobi.jacobiConjugateEquiv,
    Riemannian.Jacobi.jacobiEndpointOfVel_map_velocityPerp_eq,
    Riemannian.Jacobi.metricInner_jacobiJ_velocity_eq_zero}
  \leanok
  \uses{cor:dc-ch5-3-8, prop:dc-ch5-3-5, prop:dc-ch5-3-6, def:dc-ch5-3-1}
  Let $\gamma:[0,a]\to M$ be a geodesic in $M$, $\dim M=n$, and let
  $\mathcal{J}^\perp$ be the space of Jacobi fields with $J(0)=0$,
  $J'(0)\perp\gamma'(0)$. Let $\{J_1,\dots,J_{n-1}\}$ be a basis of
  $\mathcal{J}^\perp$. If $\gamma(t)$, $t\in(0,a]$, is not conjugate to $\gamma(0)$,
  then $\{J_1(t),\dots,J_{n-1}(t)\}$ is a basis for the orthogonal complement
  $\{\gamma'(t)\}^\perp\subset T_{\gamma(t)}M$ of $\gamma'(t)$.

  Formalized in the endpoint form ($t=a$), for a geodesic whose image lies in one chart
  source.  Under the initial-velocity parametrization ($w=J'(0)$) the space $\mathcal J^\perp$
  is $W=\ker(\text{velocityFunctional}\,g\,\gamma(0)\,\gamma'(0))$ and the target
  $\gamma'(a)^\perp$ is $WL=\ker(\text{velocityFunctional}\,g\,\gamma(a)\,\gamma'(a))$
  (\cref{cor:dc-ch5-3-8}, dimension $n-1$ each).  The heart is the new
  \emph{intrinsic moving-base pairing} \lstinline{metricInner_jacobiJ_velocity_eq_zero}: the
  fixed-chart affine law (\cref{prop:dc-ch5-3-6}, \cref{cor:dc-ch5-3-8}) upgrades to
  $\langle J(t),\gamma'(t)\rangle_g=0$ at the moving foot $\gamma(t)$ by localizing the manifold
  Jacobi field to a chart (\lstinline{isJacobiFieldOn_of_mem_source}) and transferring through the
  chart$\leftrightarrow$intrinsic bridges (\lstinline{metricInner_eq_chartMetricInner_rep},
  \lstinline{chartVectorRep_velocity}); hence the endpoint map $\Theta:J'(0)\mapsto J(a)$
  (\cref{prop:dc-ch5-3-9}) maps $W$ into $WL$ (\lstinline{jacobiEndpointOfVel_mem_velocityPerp}).
  Since $\gamma(a)$ is not conjugate to $\gamma(0)$, $\Theta$ is injective (\cref{prop:dc-ch5-3-5})
  and both hyperplanes have dimension $n-1$, so $\Theta$ restricts to a linear isomorphism
  $W\xrightarrow{\sim}WL$ (\lstinline{jacobiConjugateEquiv}, from the image
  $\Theta(W)=WL$ of \lstinline{jacobiEndpointOfVel_map_velocityPerp_eq}); consequently a basis of
  $\mathcal J^\perp$ maps to a basis of $\gamma'(a)^\perp$ whose vectors are the endpoint values
  $J_i(a)$ (\lstinline{jacobiConjugateBasis}).
\end{corollary}
